\section{Uniform oscillation at critical and residence cutoffs}
\label{sec:v90-prefix}
The original exceptional source is not a common-survivor restriction.
On a maximal segment its residence selector is constant, while its
initial-gradient selector is an initial interval because $z$ is strictly
increasing with roof. We prove an estimate for every Borel initial
interval of every physical segment, with no dependence on the position
of its new endpoint.

Let $c_\sigma\in[a_\sigma,b_\sigma]$ be Borel; the choices
$c_\sigma=a_\sigma$ and $c_\sigma=b_\sigma$ mean, respectively,
no segment and the whole segment. Define
\begin{equation}\label{eq:v90-prefix-source}
 f_c^\varepsilon(t)=\int f_\sigma^\varepsilon(t)
                                  \one_{\{t<c_\sigma\}}\,d\sigma.
\end{equation}
No regularity of $c_\sigma$ in the angular coordinate is assumed.
There is at most one new endpoint on each original segment.

\begin{lemma}[Crossing one interior endpoint]
\label{lem:v90-prefix-crossing}
For $0<h\le\rho/4$ and $0<\rho<\rho_0$, put
\[
 C_{c,h}^\varepsilon(J)=\int_J\int
    f_\sigma^\varepsilon(t)\one_{\{t<c_\sigma\le t+h\}}\,d\sigma\,dt.
\]
Then
\begin{align}
 C_{c,h}^\varepsilon(J)\le{}&\int_J e_\rho^-(t)\,dt
       +C\frac h\rho\int_{J+(-\rho,h)}p(u)\,du\notag\\
       &+C h\rho^{-3/8}\int_{J+(-\rho,h)}v(u)\,du.
                                           \label{eq:v90-prefix-crossing}
\end{align}
The constants do not depend on $c$, $\varepsilon$, the collision
count or the number of physical segments.
\end{lemma}
\begin{proof}
If $c_\sigma-a_\sigma<\rho$, every contributing starting roof
$t$ has $0<t-a_\sigma<\rho$. Its original weighted mass is bounded
by the first term. This includes endpoints arbitrarily close to a
critical entrance; no trace there is assumed.

For $c_\sigma-a_\sigma\ge\rho$, choose an anchor
$u_0\in(c_\sigma-\rho/2,c_\sigma-\rho/4)$. Every contributing
$t\in(c_\sigma-h,c_\sigma)$ lies after the anchor. On the entire
interval between them, $u-a_\sigma\ge\rho/2$, so that
$z(u)^2\ge\lambda_\circ\rho$. The exact angular-density derivative
therefore yields
\[
 A_\sigma(t)\le A_\sigma(u_0)
    +C\rho^{-3/8}\int_{c_\sigma-\rho/2}^{c_\sigma}
                             A_\sigma(u)I_m(u)z(u)^{-1/4}\,du.
\]
Average over the anchor interval of length $\rho/4$. Use
$W^\varepsilon\le1$ and integrate over a subset of the crossing
interval of length at most $h$. The resulting two bounds are
$4h/\rho$ times the segment's unweighted anchor mass and
$Ch\rho^{-3/8}$ times its fractional source mass.

For every contributing segment $c_\sigma\in J+(0,h)$, so its
anchor and comparison intervals lie in $J+(-\rho,h)$. Summing one
such interval per original segment counts each interior physical point
at most once under the Borel bijection. Thus the sums are submeasures
of the original $p$ and $v$ sources on that roof interval. This proves
\eqref{eq:v90-prefix-crossing}. Compact exhaustion and Tonelli remove
any endpoint truncations. The proof has used values inside the segment,
not a postulated trace on the potentially irregular graph $c_\sigma$.
\end{proof}

\begin{theorem}[A cutoff-uniform microscopic oscillation envelope]
\label{thm:v90-prefix-envelope}
For every fixed $L$, every one fixed admissible band and all sufficiently
small $h>0$,
\begin{equation}\label{eq:v90-prefix-envelope}
 \sup_{m,R,n,k,c}\sup_{|J|\le L}
 m^2\int_J|f_c^\varepsilon(t+h)-f_c^\varepsilon(t)|\,dt
       \le K_{B_*,L}\bigl(h^{1/33}+\varepsilon^{-1}\sqrt h\bigr).
\end{equation}
In particular, for every $a,\delta>0$, the actual source
\begin{equation}\label{eq:v90-original-union}
 R_{a,\delta}^\varepsilon(t)=\int W^\varepsilon(x)
               \one_{\{z(x)<a\ \mathrm{or}\ \ell(x)<\delta\}}\,d\mu_t(x)
\end{equation}
satisfies the same estimate, with constants independent of $a$ and
$\delta$. The union is not counted twice. The supremum over all finite
collision counts precedes the microscopic limit $h\downarrow0$.
\end{theorem}
\begin{proof}
For the zero-extended segment density and $h>0$, the elementary
product difference gives
\[
 |f_\sigma^\varepsilon(t+h)\one_{\{t+h<c_\sigma\}}
       -f_\sigma^\varepsilon(t)\one_{\{t<c_\sigma\}}|
 \le |f_\sigma^\varepsilon(t+h)-f_\sigma^\varepsilon(t)|
       +f_\sigma^\varepsilon(t)\one_{\{t<c_\sigma\le t+h\}}.
\]
Integrate before summing the complete source. The right side is
$\Xi_h^\varepsilon(J)+C_{c,h}^\varepsilon(J)$. Use
Theorem~\ref{thm:v90-whole-oscillation},
Lemma~\ref{lem:v90-prefix-crossing} and the fixed-band finite-count
envelopes \eqref{eq:v90-uniform-end}--\eqref{eq:v90-volume-finite}.
They bound it after normalization by
\[
 K_{B_*,L}\left(h^{1/32}+\varepsilon^{-1}\sqrt h
                       +\rho^{1/32}+h/\rho+h\rho^{-3/8}\right).
\]
Choose $\rho=4h^{32/33}$. For sufficiently small fixed $h$ this
lies below $\rho_0$ and obeys $h\le\rho/4$. The three new powers
are respectively $h^{1/33}$, $h^{1/33}$ and $h^{7/11}$, up to
fixed factors. This proves \eqref{eq:v90-prefix-envelope}, uniformly
in every Borel choice of endpoints. Negative shifts follow as before.

For \eqref{eq:v90-original-union}, take the full segment when
$\ell_\sigma<\delta$. On every other segment take its initial
interval on which $z<a$. Radial growth makes this interval empty,
full or a prefix with a unique threshold roof. Its endpoint is Borel:
one may determine it from the values on the countable rational roofs
of the existing segment parametrization. Thus it is a permitted
$c_\sigma$. This represents the original union exactly, not an upper
comparison to a transported or pooled source. No derivative in angle
of this new graph, and no positive lower gradient or residence cutoff,
was used in the uniform estimate.
\end{proof}

The estimate is a roof-translation modulus integrated over the starting
roof. It is stronger than knowing only the mass of the selected union,
and its constants remain uniform as both selectors tend to zero before
any collision limit. It is not an essential-supremum modulus: the
integral over $J$ in \eqref{eq:v90-prefix-envelope} must not be replaced
by a supremum over its starting roofs.
